\subsection{Discussion}

The residue-49 calculations confirm the expected catalytic advantage of selenium: Sec lowers the proton-transfer barrier relative to Cys in both human and mouse GPX6.\cite{Reich2016,prabhakar_elucidation_2005,cardey_selenocysteine_2007,rees_ancient_2024,Kanavos2026SecSelectivity} However, the effect is not the same in the two proteins. Replacing Sec with Cys raises the barrier by 1.43~kcal/mol in human GPX6, whereas replacing Cys with Sec lowers it by 2.84~kcal/mol in mouse GPX6. The residues surrounding the active site therefore modulate the catalytic effect of Sec.\cite{Warshel2006,Kamerlin2010,VillaWarshel2001}

The mutational calculations also show epistasis. Human substitutions added to mouse Sec-GPX6 generally lower the barrier or leave it unchanged, whereas several mouse substitutions added to human Cys-GPX6 raise the barrier. The same sequence difference can therefore have a different effect in the two proteins.\cite{Weinreich2006,starr2016epistasis,storz_compensatory_2018,Kaltenbach2015} In some cases, the effect even changes sign as other substitutions accumulate. The number of local hydrogen bonds does not explain these changes, which indicates that several residues contribute jointly to transition-state stabilization. The selected paths are calculated low-barrier routes, not reconstructions of the historical mutation order. They do not include population genetics, changes outside the 20 selected residues, or evolutionary steps that may have been favorable for functions other than the reaction studied here.

Rees et al. inferred five independent losses of Sec in mammalian GPX6 and found higher evolutionary rates and convergent changes in the corresponding branches, particularly in the catalytic domain.\cite{rees_ancient_2024} They interpreted the Cys-containing lineages as convergent adaptations shaped by pleiotropy and epistasis, rather than as simple reversions to an ancestral state. Our calculations provide an energetic basis for this interpretation. Human GPX6 depends more strongly on Sec, whereas mouse GPX6 has accumulated changes that support catalysis with Cys and can also produce a low barrier when Sec is reintroduced. The two calculated paths overlap in the triangular sequence projection, so the results do not support strict hysteresis. Nevertheless, the order favored by the EVB calculations is not random: several substitutions recur early in the lower-barrier paths. GPX6 evolution therefore appears constrained, but not irreversible.

The charge model is an important limitation. Maestro assigned the same resonance charges to sulfur in cysteine and selenium in selenocysteine. This does not mean that the two atoms are electronically equivalent; it reflects the parameterization procedure and the more limited force-field treatment of selenium. Sulfur and selenium differ in size, polarizability, and electronegativity, and charges derived from dedicated quantum calculations or experimental calibration could change the numerical Sec--Cys differences. However, all variants were calculated with the same parameters, and the Sec/Cys exchange still had different effects in human and mouse GPX6.

The transferring proton was treated classically in the EVB simulations. The diabatic potentials were calibrated against a quantum-mechanical free-energy profile, but the calculations did not include zero-point energy or proton tunneling and therefore cannot predict kinetic isotope effects. These contributions can be substantial in enzymatic proton transfer and may change the absolute barriers.\cite{VillaWarshel2001} They could be estimated with a separate quantum correction or a path-integral calculation. Because Sec and Cys variants were treated with the same mechanism, calibration, and simulation protocol, much of the common quantum correction is expected to cancel in their comparison. The cancellation may be incomplete if the variants change donor--acceptor distances or force constants. Explicit calculations would be needed to determine its magnitude.

The ProteinMPNN analysis is exploratory. It tests which sequences are compatible with the sampled structures, whereas the EVB calculations rank substitutions by their catalytic barriers. ProteinMPNN recovered little of the node~25 sequence, showing that structural compatibility alone is not sufficient to reconstruct the ancestor. The structural inputs are also asymmetric: mouse GPX6 is represented by an experimental structure (PDB 7FC2), whereas human GPX6 is an AlphaFold model. This difference may favor mouse-like sequences. Future comparisons should use either experimental structures for both proteins or models generated in the same way.

The present calculations cover 20 positions near the active site, so interactions involving more distant residues remain unexplored. Larger variant sets could test whether these interactions alter the preferred paths or improve recovery of the ancestral sequence. For enzyme design, sequence models could be used to propose variants and EVB calculations to rank them by catalytic barrier, including tests in more than one protein background.\cite{Tran2026Science}
